El error cuadrático medio es
\[
E_N=\frac8{\pi^2}\sum_{j=N}^{\infty}\frac1{(2j+1)^2}
\le\frac4{\pi^2(2N-1)}.
\]
Por tanto, \(N=21\) términos garantizan \(E_N\le0.01\). En cambio, \(S_N(0)=0\) y \(f(0)=1\), así que el error puntual en cero siempre es uno. Reasignar \(f(0)=0\) anula ese error aislado sin alterar coeficientes ni \(E_N\), pero no produce un error uniformemente pequeño cerca del salto: en cualquier entorno de cero, el supremo del error sigue siendo al menos uno.
